\section{Introduction}\label{sec:introduction}

Scientific communication still centers on the paper, even as research increasingly depends on software, data, computational environments, formal specifications, and machine-assisted analysis. A manuscript remains effective for presenting an argument to human readers, establishing priority, and placing a result in its disciplinary context. It is less effective as an interface for reconstructing the exact chain of claims on which a conclusion depends. A reader who wants to verify or reuse one result must often recover the relevant source version, identify the precise statement imported from each citation, reconcile notation and assumptions across papers, locate executable artifacts, and determine which parts of the argument have actually been checked. These tasks become more difficult as literature indexes, repositories, software packages, and model outputs change over time.

Several lines of work address parts of this problem. Research Objects treat a scholarly result as more than a manuscript by packaging data, methods, workflows, and provenance into a shareable unit~\cite{bechhofer2013linked}. Workflow-centric Research Objects refine this approach with ontologies that preserve the structure and evolution of computational analyses~\cite{belhajjame2015ontologies}. RO-Crate provides a practical linked-data format for packaging research artifacts and their contextual metadata~\cite{soilandreyes2022rocrate,rocrate2024spec}. The FAIR principles establish that research data should be findable, accessible, interoperable, and reusable, with particular emphasis on machine actionability~\cite{wilkinson2016fair}. PROV-O supplies a general vocabulary for entities, activities, agents, and provenance relations~\cite{w3c2013provo}, while the Common Workflow Language supports portable and reusable computational workflows~\cite{crusoe2022cwl}. Together, these efforts provide important infrastructure for identifying artifacts, describing their origins, and preserving executable processes.

Artifact packaging and provenance do not, by themselves, resolve the scientific dependency structure within those artifacts. A well-formed research package may identify a manuscript, dataset, workflow, and software environment without stating which source span supports a particular claim, which earlier theorem is load-bearing for that claim, or whether the imported theorem is used under compatible assumptions. Provenance can show that an output was generated by an activity using specified inputs; it need not show that a scientific conclusion follows from those inputs. This distinction does not weaken the value of Research Objects, RO-Crate, FAIR, PROV-O, or workflow standards. It identifies a complementary layer that a verification-aware system must add.

A second relevant line of work concerns commitments made before outcomes are known. Preregistration records hypotheses, designs, and analysis plans before data analysis, thereby separating confirmatory decisions from choices made after observing results~\cite{nosek2018preregistration}. Registered Reports incorporate this principle into publication by reviewing the question and method before the results are available~\cite{chambers2013registered}. These practices were developed primarily to improve study design and evidential interpretation, not to freeze the state of a literature search. Their central insight nevertheless applies more broadly: a research process is easier to audit when its initial question, decision rules, and subsequent deviations are recorded before the outcome is selected.

AgtXIv synthesizes this insight with provenance and versioned publication through a design object called a \emph{Scientific Query Freeze}. This term denotes an AgtXIv protocol element rather than an established concept from preregistration, publication protocols, or claim-tracing systems. A Scientific Query Freeze is a versioned, timestamped record of the research question, exact retrieval queries, searched sources or endpoints, temporal cutoff, filters, inclusion and exclusion rules, retrieved candidate identifiers or a hash of the candidate set, software and model versions, and documented deviations from the initial procedure. Freezing only a query string is insufficient because a literature index can return different records for the same string after additions, deletions, deduplication, or ranking changes. Recording the candidate set, or a stable digest together with a recoverable manifest, separates changes in the index from changes in the investigator's selection process. The resulting object does not guarantee an unbiased or complete review. It establishes an auditable acquisition boundary that can be connected to preregistered decisions, provenance records, and immutable publication releases.

A third line of work moves from artifacts to relations among scholarly statements. CiTO provides typed citation relations that distinguish different functions of a citation rather than treating every bibliographic link as equivalent~\cite{shotton2010cito}. The scite citation index similarly exposes citation context and classifies how later work refers to an earlier publication~\cite{nicholson2021scite}. ClaimFlow proposes a claim-centered account of scientific development with five directed relation types: background, support, extension, qualification, and refutation~\cite{pramanick2026claimflow}. These relations describe how one locally grounded claim interacts with another. In particular, a citation-grounded support relation states that one claim supplies support in a specified context; it does not establish that the supported claim is globally true, that all of its assumptions hold, or that no contrary evidence exists.

This observation requires several relation families to remain distinct. A bibliographic relation connects documents or bibliographic records. An epistemic relation describes how one claim supports, extends, qualifies, refutes, or provides background for another. A mathematical dependency relation states that a definition, assumption, lemma, theorem, or inference step is required to derive a selected mathematical conclusion. A provenance or evidence relation records where a statement, artifact, observation, or verification result came from and what it directly supports. These relations can inform one another, but none can be substituted automatically for another. A citation is not necessarily a proof dependency; a proof dependency is not necessarily empirical support; and a provenance record is not a correctness certificate.

Scientific claim verification provides another component of the required infrastructure. SciFact frames verification as deciding whether scientific evidence supports or refutes a claim and identifying the relevant evidence spans~\cite{wadden2020scifact}. This formulation demonstrates the value of claim-level evidence retrieval, but general scientific verification requires more than one entailment label. Mathematical correctness, source fidelity, physical interpretation, empirical support, and computational reproducibility involve different objects and checking procedures. Nanopublications address a related granularity problem by packaging a minimal assertion with provenance and publication information~\cite{kuhn2016nanopublications}. They show how small, provenance-aware assertions can become independently identifiable publication units. AgtXIv adopts the need for minimal claim objects while adding typed dependencies, query-relative closure, and heterogeneous verification records.

Recent proposals make publications directly accessible to agents. The Agentic Publication Protocol defines a lightweight, version-controlled repository format that combines a canonical paper with code, data, environment information, reproducibility instructions, and agent guidance~\cite{lu2026app}. Agent-Native Research Artifacts propose a richer machine-readable research object in which scientific logic, executable components, evidence, and research traces can be represented directly~\cite{liu2026ara}. Paper2Agent converts papers and associated code into interactive agents that expose methods and resources through executable interfaces~\cite{miao2025paper2agent}. These projects occupy different operating points. The Agentic Publication Protocol emphasizes author-side packaging and release; Agent-Native Research Artifacts propose a more deeply structured artifact; Paper2Agent emphasizes executable access to a paper's methods. These proposals suggest that agents can consume, explain, and operate scientific artifacts, but they do not make every agent response a verified scientific claim. A paper agent may faithfully expose an executable method while leaving the source-to-claim dependency chain implicit, and an executable result may remain scientifically misaligned with the claim for which it is used.

The unresolved problem is therefore not simply how to package more material with a paper or how to place a conversational interface over it. Scientific discovery needs infrastructure that can answer a narrower and more demanding question: for a selected target claim, which exact source statements and dependencies are required, which parts have passed which checks, and what unresolved local work remains? Answering this question requires connecting publication identity, claim structure, and verification without collapsing them into one representation or one trust score.

The methodological thesis of AgtXIv is that agent-native scientific publishing and discovery require three connected but non-collapsed layers. The first layer consists of immutable, versioned publication and source objects, including frozen manuscripts, repositories, manifests, hashes, and Scientific Query Freezes. The second consists of claim-level epistemic and dependency graphs grounded in exact source spans. The third consists of verification records that separately track source fidelity, formal mathematical correctness, formal and semantic alignment, empirical support, and computational reproducibility. Connections among these layers make a claim traceable and reusable. Their separation prevents success on one axis from being reported as success on all others.

AgtXIv is a math-first, verification-aware research infrastructure and experimental protocol built around this thesis. A frozen paper is represented by a source-bounded \emph{PaperAgent}, but the principal reusable unit is a claim-level \emph{MathContract}. A MathContract records a normalized mathematical statement, its assumptions, imports, source manifestations, formal declarations when available, alignment results, version, and unresolved blockers. PaperAgents organize source-bounded contributions and interactions; they are not trusted by authority, and freezing a paper does not turn its claims into axioms.

The system maintains different graph views for different purposes. A paper-level interaction graph records typed bibliographic and epistemic relations and may contain cycles. Companion papers can use different claims from one another without forming a circular proof. For a selected query $q$, AgtXIv derives a load-bearing mathematical dependency graph $G_{\mathrm{math}}(q)$. Its nodes represent definitions, assumptions, lemmas, theorems, constructions, explicit foundations, and inference steps required by the target. After duplicate or equivalent claims have been resolved, this query-relative mathematical view must be acyclic so that it defines a valid dependency and build order. A cycle at this level indicates a circular import, an incorrectly oriented edge, an unresolved equivalence, or a confusion between epistemic interaction and proof dependence. When an inference requires several premises jointly, AgtXIv represents the inference as an explicit node: all premises point to the inference step, which points to its conclusion. This construction preserves multi-premise semantics while using an ordinary typed directed graph.

Candidate construction remains deliberately separate from acceptance. Literature retrieval, large language model extraction, theorem search, symbolic algebra, and related tools act as \emph{oracles}: they may propose claims, source anchors, relations, decompositions, formal declarations, proofs, or computational certificates. Their outputs remain candidates until the configured checks succeed. Model confidence, citation counts, and retrieval rank do not establish scientific truth. This policy permits extensive automation while preventing proposal mechanisms from silently promoting their own outputs.

For mathematical subclaims, AgtXIv uses Lean~4 and pinned formal packages to check whether a conclusion follows from the encoded assumptions. The Lean kernel provides a precise mathematical correctness check for the formal declaration; it does not show that the declaration faithfully represents the source. AgtXIv therefore pairs formalization with source-blind backtranslation and source-aware alignment auditing. The backtranslator receives the formal declaration and the definitions needed to interpret it, but not the original source claim, and reconstructs the mathematical statement that the declaration actually expresses. An alignment auditor then compares the frozen source, normalized claim, formal declaration, and backtranslation for changes in quantifiers, object types, assumptions, domains, exactness, approximation status, and conclusion strength. This procedure cannot eliminate semantic error, but it makes common forms of formalization drift inspectable.

The distinction between correctness and alignment is especially important in theoretical physics and other mathematically structured sciences. Lean can prove that a theorem follows under encoded assumptions. It cannot, by that fact alone, prove that the formal object represents the intended physical system, that an approximation is valid in the claimed regime, that an observable matches an experimental procedure, that empirical evidence supports the model, or that numerical software reproduces a reported result. AgtXIv therefore records these dimensions separately. It does not assign one scalar trust score to a paper or infer truth from the number or direction of citations.

Verification proceeds over a query-relative dependency closure rather than an entire field. The system traces a target backward to accepted reusable contracts, query-relative root papers, explicit external foundations, or a visible unresolved frontier. It then builds forward, importing accepted contracts and checking only the residual local mathematical delta. Failed proof construction, missing sources, or alignment errors trigger local graph repair rather than unrecorded narrative revision. The intended result is a reusable verification receipt containing the target claim, source addresses, accepted and conditional imports, checked declarations, alignment audits, dependency versions, and the first unresolved frontier.

The intended contributions are consequently framed as a protocol, architecture, and pilot agenda rather than a completed universal platform:
\begin{enumerate}[leftmargin=*]
    \item AgtXIv defines a frozen acquisition boundary that combines Scientific Query Freezes with immutable, versioned source and publication objects.
    \item It separates typed literature interactions from query-relative, load-bearing mathematical dependencies while preserving source-span grounding and explicit multi-premise inference steps.
    \item It specifies claim contracts and reusable verification records with orthogonal status axes for provenance, mathematics, alignment, semantics, evidence, and computation.
    \item It proposes an incremental formalization loop that reuses accepted root contracts, checks local deltas in pinned Lean environments, and exposes blocked or disputed dependencies instead of hiding them.
\end{enumerate}

The immediate scope is intentionally narrow: one bounded dependency closure in theoretical physics or another mathematically structured science, with a small number of source papers and at least one load-bearing formal reconstruction. AgtXIv does not attempt to formalize complete papers, automate field-wide review, replace expert physical judgment, verify arbitrary numerical software, or establish universal scientific truth. Its pilot question is whether a source-grounded, claim-level, verification-aware build process can make one scientific dependency path more auditable and reusable than a paper list, narrative summary, or executable agent alone.

The remainder of the paper will develop the graph and system architecture, specify the claim contracts and verification semantics, describe the incremental pilot workflow, and discuss the scope, limitations, and evaluation criteria.
